% =============================================================================
% bubble_sheet_body.tex --- shared drawing logic for the bubble-sheet
% templates. Not compiled directly; \input from bubble_sheet.tex (50q) and
% bubble_sheet_100q.tex (100q).
%
% Defaults below are picked up unless the wrapper has \newcommand'd them
% first. Counts (\NumQuestions / \NumChoices / \NumColumns) have no
% defaults because compiling without choosing them is almost certainly a
% mistake -- the wrapper must set them.
% =============================================================================

\providecommand{\BubbleRadius}{2.0}      % drawn circle radius, mm
\providecommand{\AnchorRadius}{1.85}     % column anchor disc radius, mm
\providecommand{\NumberFont}{\small}     % question-number font
\providecommand{\NumberOffset}{4}        % mm from anchor centre to number's right edge
\providecommand{\ChoiceFont}{\footnotesize\bfseries}
\providecommand{\QStartOffset}{0}        % added to \q on every row (multi-page sheets)
\providecommand{\PageLabel}{}            % shown right of header subtitle if non-empty

\begin{tikzpicture}[x=1mm, y=-1mm, line width=0.4pt]
  \useasboundingbox (0,0) rectangle (210, 297);

  % -- Derived quantities -------------------------------------------------
  % Distribute NumQuestions across NumColumns. The first ExtraRows columns
  % carry one extra question each, matching web/src/types.ts :: defaultColumns.
  \pgfmathtruncatemacro{\BaseRows}{int(\NumQuestions / \NumColumns)}
  \pgfmathtruncatemacro{\ExtraRows}{\NumQuestions - \BaseRows * \NumColumns}
  \pgfmathtruncatemacro{\NumChoicesM}{\NumChoices - 1}

  % Bubble pitch and row height. Must match BUBBLE_X_PITCH_RATIO and
  % ROW_HEIGHT_RATIO in web/src/cv/grid.ts (both expressed as fractions of
  % the 180mm hSpan): 6/180 and 7/180.
  \pgfmathsetmacro{\XPitch}{6.0}
  \pgfmathsetmacro{\RowHeight}{7.0}

  % Column anchors span x = 30mm to x = (183 - last-bubble offset), leaving a
  % ~12mm quiet zone to the right-corner marker at x=195.
  \pgfmathsetmacro{\AnchorXFirst}{30}
  \pgfmathsetmacro{\AnchorXLast}{183 - \NumChoicesM * \XPitch}
  \pgfmathsetmacro{\ColSpacing}{%
    \NumColumns > 1 ? (\AnchorXLast - \AnchorXFirst) / (\NumColumns - 1) : 0}

  % Row-0 y = anchor y + one row height (matches computeGridParams in grid.ts).
  \pgfmathsetmacro{\AnchorY}{60}
  \pgfmathsetmacro{\RowOneY}{\AnchorY + \RowHeight}

  % Tallest column drives the bottom-fiducial y. Mirrors the maxRows term in
  % web/src/cv/rectify.ts :: templateHeightMm so the scanner can rectify any
  % size sheet without a per-template constant.
  \pgfmathtruncatemacro{\MaxRows}{\BaseRows + (\ExtraRows > 0 ? 1 : 0)}
  \pgfmathsetmacro{\BottomCornerY}{\AnchorY + \MaxRows * \RowHeight + 7.17}

  % -- Page header -------------------------------------------------------
  \node[anchor=north west, font=\LARGE\bfseries] at (15, 16)
        {Exam Answer Sheet};
  \node[anchor=north west, font=\footnotesize\itshape, text=black!55,
        text width=140mm] at (15, 25)
        {Fill each bubble completely. To change an answer, erase fully or
         strike through the wrong mark and fill the new one.};
  \node[anchor=north east, font=\footnotesize\itshape, text=black!55]
        at (195, 25) {\PageLabel};

  \node[anchor=base west, font=\small] at (15, 44)  {Name};
  \draw (28, 44) -- (130, 44);
  \node[anchor=base west, font=\small] at (140, 44) {Date};
  \draw (152, 44) -- (195, 44);

  \draw[draw=black!30, line width=0.3pt] (15, 49) -- (195, 49);

  % -- Fiducial corner squares -------------------------------------------
  % TL and TR define the hSpan that the scanner uses to calibrate the grid.
  % BL and BR help confirm orientation; one may be obscured without breaking
  % detection. Bottom y comes from \BottomCornerY so the layout adapts to
  % \NumQuestions automatically (215mm at 50q, ~285mm at 100q).
  \foreach \cx in {15, 195} {%
     \fill (\cx-1.75, 55-1.75) rectangle (\cx+1.75, 55+1.75);
     \fill (\cx-1.75, \BottomCornerY-1.75) rectangle (\cx+1.75, \BottomCornerY+1.75);
  }

  % -- Columns -----------------------------------------------------------
  \foreach \col in {1,...,\NumColumns} {%
     \pgfmathsetmacro{\ax}{\AnchorXFirst + (\col - 1) * \ColSpacing}
     \pgfmathtruncatemacro{\rowsCol}%
        {\BaseRows + (\col <= \ExtraRows ? 1 : 0)}
     \pgfmathtruncatemacro{\colMinus}{\col - 1}
     \pgfmathtruncatemacro{\qStart}%
        {1 + \colMinus * \BaseRows
           + (\colMinus < \ExtraRows ? \colMinus : \ExtraRows)}

     % Column anchor disc (solid black). Sits 5mm below the top corner row
     % so its centre is within ~40 px of the corner-y average --- inside the
     % scanner's 50-px window for "same header row".
     \fill (\ax, \AnchorY) circle[radius=\AnchorRadius mm];

     % Choice-letter labels between the anchor and the first bubble row.
     % Offset 2.5mm fits the 7mm row height: there's only ~3mm of clear space
     % between anchor bottom and row-1 bubble top, and \footnotesize letters
     % need most of it.
     \foreach \ci in {0,...,\NumChoicesM} {%
        \pgfmathsetmacro{\bx}{\ax + \ci * \XPitch}
        \pgfmathtruncatemacro{\code}{97 + \ci}  % ASCII 'a' + ci
        \node[anchor=south, font=\ChoiceFont, inner sep=0.5pt]
              at (\bx, \RowOneY - 2.5) {\char\code};
     }

     % Question rows. Row 1 centre = anchor_y + row_height (matches
     % web/src/cv/grid.ts :: computeGridParams).
     \foreach \r in {1,...,\rowsCol} {%
        \pgfmathsetmacro{\yy}{\RowOneY + (\r - 1) * \RowHeight}
        \pgfmathtruncatemacro{\q}{\qStart + \r - 1 + \QStartOffset}

        \node[anchor=east, font=\NumberFont] at (\ax - \NumberOffset, \yy) {\q};

        \foreach \ci in {0,...,\NumChoicesM} {%
           \pgfmathsetmacro{\bx}{\ax + \ci * \XPitch}
           \draw (\bx, \yy) circle[radius=\BubbleRadius mm];
        }
     }
  }
\end{tikzpicture}
